%%%%%%%%%%%%%%%%%%%%%%%%%%%%%%%%%%%%%%%%%%%%%%%%%%%%%%%%%%%%%%%%%%%%%%
% Lớp 11 | Đường thẳng, mặt phẳng. Quan hệ song song trong | Đường thẳng và mặt phẳng song song
% Dạng 1H4?3-1: Câu hỏi lý thuyết
% Số câu: 54  (sinh tự động bởi tach_cau.py; từ nay thầy sửa trực tiếp file này)
%%%%%%%%%%%%%%%%%%%%%%%%%%%%%%%%%%%%%%%%%%%%%%%%%%%%%%%%%%%%%%%%%%%%%%

%%%%%%%%%%%%%%%% TRẮC NGHIỆM 4 PHƯƠNG ÁN %%%%%%%%%%%%%%%%

% nguồn: Ok_Toan11_TracNghiem_Lan2
\begin{ex}%[1H4N3-1]%[ma:c54b8d1c]
	Cho hai đường thẳng phân biệt $ a$, $b$ và mặt phẳng $( \alpha)$. Giả sử $ a \parallel ( \alpha)$ và $ b \parallel ( \alpha)$. Mệnh đề nào sau  đây đúng?
	 \choice
		{ $ a$ và $ b$ không có điểm  chung}
		 { $ a$ và $ b$ hoặc song song hoặc chéo nhau}
		{\True  $ a$ và $ b$ hoặc song song hoặc chéo nhau hoặc cắt nhau}
		{ $ a$  và $ b$ chéo nhau}
	\loigiai{ 
		Ta có $ a$ và $ b$  hoặc song song hoặc chéo nhau hoặc cắt nhau.
	 }
\end{ex}

% nguồn: Ok_Toan11_TracNghiem_Lan2
\begin{ex}%[1H4N3-1]%[ma:30e25871]
	Cho đường thẳng $ a$ nằm trong mặt phẳng $( \alpha)$. Giả sử  $ b\not{ \subset }( \alpha)$. Mệnh đề nào sau đây đúng?
	 \choice
		{Nếu $ b \parallel ( \alpha)$ thì $  b \parallel a $}
		{Nếu $ b$ cắt $( \alpha)$ thì $ b$ cắt  $ a $}
		{\True Nếu $ b \parallel  a$ thì $ b \parallel ( \alpha) $}
		{Nếu $ b$ cắt $( \alpha)$ và $( \beta)$ chứa $ b$ thì giao tuyến của $( \alpha)$ và $( \beta)$ là đường  thẳng cắt cả $ a$ và $ b $}
	 \loigiai{
		\begin{itemize}
			\item Khẳng định "Nếu $  b \parallel ( \alpha)$ thì $ b \parallel a $" sai. Vì nếu $ b \parallel ( \alpha)$ thì $ b \parallel a$ hoặc $ a$, $b$ chéo nhau.
			\item Khẳng định "Nếu $ b$ cắt $( \alpha)$ thì $ b$  cắt $ a $" sai. Vì nếu $ b$ cắt $( \alpha)$ thì $ b$ cắt $ a$ hoặc $ a$, $b$ chéo nhau.
			\item Khẳng định "Nếu $ b$ cắt $( \alpha)$ và $( \beta)$ chứa $ b$ thì giao tuyến của $( \alpha)$ và $( \beta)$ là đường thẳng cắt cả $ a$ và $ b $" sai. Vì nếu $ b$ cắt $( \alpha)$ và $( \beta)$ chứa $ b$ thì giao tuyến của $( \alpha)$ và $( \beta)$ là đường thẳng cắt $ a$ hoặc  song song với $ a$.
		\end{itemize}
	 }
\end{ex}

% nguồn: Ok_Toan11_TracNghiem_Lan2
\begin{ex}%[1H4N3-1]%[ma:26f328a0]
	Cho $ d\parallel ( \alpha)$, mặt phẳng $( \beta)$ qua $ d$ cắt $(  \alpha)$ theo giao tuyến $d'$. Khẳng định nào sau đây đúng?
	 \choice
		{\True  $ d\parallel d' $}
		{$  d$ cắt $d'$}
		{$ d$  và $d'$ chéo nhau}
		 {$ d\equiv d' $}
	\loigiai{ 
		Ta có  $d'=( \alpha) \cap ( \beta)$. Do $ d$ và $d'$ cùng thuộc $( \beta)$ nên $ d$ cắt $d'$ hoặc $ d \parallel d'$ hoặc $d\equiv d'$.
		\begin{itemize}
			\item Nếu $d\equiv d' \Rightarrow  d \subset \left(\alpha\right)$ (vô lí).
			\item Nếu $ d$ cắt $d'$ suy ra $d$ cắt $(\alpha)$ (vô lí,  do $d\parallel (\alpha)$).
		\end{itemize}
		 Vậy $ d \parallel d'$.
	} 
\end{ex}

% nguồn: Ok_Toan11_TracNghiem_Lan2
\begin{ex}%[1H4N3-1]%[ma:05ebd93a]
	Cho hai đường thẳng chéo nhau  $ a$ và $ b$. Khẳng định nào sau đây \textbf{sai}?
	 \choice
		{\True Có duy nhất một mặt phẳng song song với $ a$ và  $ b $}
		{Có duy nhất một mặt phẳng qua $ a$ và song song  với $ b $}
		{Có duy nhất một mặt phẳng qua điểm $M$, song song  với $ a$ và $ b$}
		{Có vô số đường thẳng song song với $ a$ và  cắt $ b $}
	\loigiai{ 
		Có có vô số mặt phẳng song song với $2$ đường thẳng chéo  nhau.
	} 
\end{ex}

% nguồn: Ok_Toan11_TracNghiem_Lan2
\begin{ex}%[1H4N3-1]%[ma:9da17bf8]
	Trong không  gian, khẳng định nào sau đây là khẳng định đúng?
	\choice 
		{Nếu đường thẳng $a$ song song với đường  thẳng $b$ nằm trong mặt phẳng $\left(P\right)$ thì $a$ song song với $\left(P\right)$}
		{\True  Cho đường thẳng $a$ song song với mặt phẳng $\left(P\right)$. Nếu mặt phẳng $\left(Q\right)$ chứa $a$ và cắt $\left(P\right)$ theo giao tuyến $d$ thì $d$ song song với $a$}
		{Cho hai mặt phẳng phân biệt $\left(P\right)$, $\left(Q\right)$  lần lượt chứa hai đường thẳng song song $a,b$. Nếu $\left(P\right)$ và $\left(Q\right)$ cắt nhau theo giao tuyến $d$ thì $d$ song song với $a$}
		{Nếu hai đường thẳng phân biệt $a,b$  cùng song song với mặt phẳng $\left(P\right)$ thì $a$, $b$ song song nhau}
	 \loigiai{
		\begin{itemize}
			\item Nếu đường thẳng $a$ song song với đường thẳng $b$ nằm trong mặt phẳng $\left(P\right)$ thì  $a$ có thể song song với $\left(P\right)$ hoặc $a\subset \left(P\right)$.
			\item Cho đường thẳng $a$ song song với mặt phẳng $\left(P\right)$. Nếu mặt phẳng $\left(Q\right)$ chứa $a$ và cắt $\left(P\right)$ theo giao tuyến $d$ thì $d$ song song  với $a$ là mệnh đề đúng.
			\item Cho hai mặt phẳng phân biệt $\left(P\right)$, $\left(Q\right)$ lần lượt chứa hai đường thẳng song song $a,b$. Nếu $\left(P\right)$ và $\left(Q\right)$ cắt nhau theo giao tuyến $d$ thì $d$ song song với $a$ hoặc  $d$ trùng với $a$.
			\item Nếu hai đường thẳng phân biệt $a$, $b$ cùng song song với mặt phẳng $\left(P\right)$ thì $a$, $b$ có thể cắt nhau, song song  với nhau hoặc chéo nhau.
		\end{itemize}
	} 
\end{ex}

% nguồn: Ok_Toan11_TracNghiem_Lan2
\begin{ex}%[1H4N3-1]%[ma:34c4b49d]
	Cho $ d\parallel ( \alpha)$, mặt phẳng $( \beta)$ qua $ d$ cắt $( \alpha)$ theo giao tuyến $d'$.  Khẳng định nào sau đây là đúng?
	 \choice
		{ $  d$ cắt $d'$}
		{\True $  d\parallel d'$}
		 { $ d$ và $d'$ chéo nhau}
		{ $ d\equiv  d'$}
	 \loigiai{
		Ta có ${d'}=( \alpha) \cap ~( \beta)$. Do $ d$ và $d'$ cùng thuộc $( \beta)$ nên $ d$ cắt $d'$ hoặc $  d\parallel d'$. \\
		Nếu  $ d$ cắt $d'$ thì $ d$ cắt $( \alpha)$ (mâu thuẫn với giả thiết).\\
		Vậy $ d\parallel d'$. 
	 }
\end{ex}

% nguồn: Ok_Toan11_TracNghiem_Lan2
\begin{ex}%[1H4N3-1]%[ma:b6416411]
	Cho hai  đường thẳng phân biệt $ a,\ b$ và mặt phẳng $( \alpha)$. Giả sử $ a \parallel ( \alpha)$, $ b \subset ( \alpha)$. Khi đó
	 \choice
		{$ a  \parallel b $}
		{$ a$, $b$ chéo  nhau}
		{\True $ a \parallel b$ hoặc $ a$, $b$ chéo  nhau}
		{$ a$, $b$ cắt  nhau}
	\loigiai{ 
		\begin{center}
			\begin{tikzpicture}[scale=1,line cap=round,line join=round,font=\footnotesize,>=stealth]
					\def\dai{3.5} \def\rong{2} \def\goc{70} 
					\path  (0,0) coordinate (A)++(0:\dai) coordinate (B)
						++(\goc:\rong)  coordinate (C)++(180:\dai) coordinate (D)
						($(A)+(50:1.2)$) coordinate (P)++(0:2.7) coordinate  (Q)
						($(P)+(90:1.7)$) coordinate  (P')++(0:2.7) coordinate (Q');
					\draw  (A)--(B)--(C)--(D)--cycle (P)--node[above,midway]{$b$}
						(Q) 
						 (P')--node[above,midway]{$a$}
						(Q'); 
					 \draw pic[draw,angle radius=5mm]{angle=B--A--D};
					 \draw (A) node[shift={(30:0.3)}]{$\alpha$};
				\end{tikzpicture}
			 \qquad
			\begin{tikzpicture}[scale=1,line cap=round,line join=round,font=\footnotesize,>=stealth]
					\def\dai{3.5} \def\rong{2} \def\goc{70} 
					\path (0,0) coordinate (A)++(0:\dai)  coordinate (B)
						++(\goc:\rong) coordinate (C)++(180:\dai) coordinate  (D)
						($(A)+(50:1.2)$) coordinate (P)++(0:2.7) coordinate (Q) 
						($(P)+(90:1.7)$)  coordinate (P')++(0:3.2) coordinate (Q')
						($(P)+(-50:0.8)$) coordinate (P1)++(30:2.5) coordinate  (Q1);
					\draw  (A)--(B)--(C)--(D)--cycle (P)--node[above,pos=0.9]{$c$}
						 (Q)
						 (P')--node[above,midway]{$a$}
						(Q') 
						(P1)--node[below,pos=0.2]{$b$} 
						 (Q1);
					\draw pic[draw,angle  radius=5mm]{angle=B--A--D};
					\draw (A)  node[shift={(30:0.3)}]{$\alpha$};
				\end{tikzpicture}
		\end{center}
		Vì $ a  \parallel ( \alpha)$ nên tồn tại đường thẳng $ c \subset ( \alpha)$ thỏa mãn $ a \parallel c $. Suy ra $ b$, $c$ đồng phẳng và xảy ra các trường hợp sau:
		\begin{itemize}
			\item Nếu $ b$ song song hoặc  trùng với $ c$ thì $ a \parallel b$.
			\item Nếu $ b$ cắt $ c$  thì $ b$ cắt $( \beta)\equiv ( a,c)$ nên $ a$, $b$ không đồng phẳng. Do đó $ a$, $b$ chéo nhau.
		\end{itemize}
	 }
\end{ex}

% nguồn: Ok_Toan11_TracNghiem_Lan2
\begin{ex}%[1H4N3-1]%[ma:8f849ef6]
	Cho đường thẳng $a$  và mặt phẳng $(\alpha)$. Chọn phương án đúng?
	\choice 
		{Nếu $a$ song  song với mặt phẳng $(\alpha)$ thì $a$ song song với mọi đường thẳng nằm trên mặt phẳng $(\alpha)$}
		{\True $a$ song song với mặt phẳng $(\alpha)$ khi và chỉ khi $a$ và $\left(\alpha \right)$ không có điểm  chung}
		{Nếu $a$ song song với mặt phẳng  $(\alpha)$ thì mọi mặt phẳng chứa $a$ đều song song với mặt phẳng $(\alpha)$}
		{Nếu $a$ song song với một đường thẳng song song với mặt  phẳng $(\alpha)$ thì $a$ song song với mặt phẳng $(\alpha)$}
	 \loigiai{
		Theo định nghĩa đường thẳng song song với mặt phẳng: đường thẳng $a$ song song với mặt phẳng $(\alpha)$ khi  và chỉ khi $a$ và $(\alpha)$ không có điểm chung.
	 }
\end{ex}

% nguồn: Ok_Toan11_TracNghiem_Lan2
\begin{ex}%[1H4N3-1]%[ma:67e8d8e2]
	Cho đường thẳng $a$ nằm trong mặt phẳng $\left(\alpha \right)$. Giả sử $b\not\subset \left(\alpha \right)$.  Mệnh đề nào sau đây đúng?
	 \choice
		 {Nếu $b \parallel \left(\alpha \right)$ thì $b \parallel a$}
		{Nếu $b$ cắt $\left(\alpha \right)$ thì $b$ cắt  $a$}
		{\True Nếu $b  \parallel a$ thì $b \parallel \left(\alpha \right)$}
		{Nếu $b \parallel \left(\alpha \right)$ và $\left(\beta \right)$ chứa $b$ thì $\left(\beta \right)$  sẽ cắt $\left(\alpha \right)$ theo giao tuyến là đường thẳng song song với $b$}
	\loigiai{ 
		Cho đường thẳng $a$ nằm  trong mặt phẳng $\left(\alpha \right)$. Giả sử $b\not\subset \left(\alpha \right)$.\\
		Nếu $b \parallel a$ thì $b  \parallel \left(\alpha \right)$.
	 }
\end{ex}

% nguồn: Ok_Toan11_TracNghiem_Lan2
\begin{ex}%[1H4N3-1]%[ma:694e7537]
	Cho mặt phẳng $( P)$ và hai đường thẳng song song $ a$ và $ b$. Khẳng định  nào sau đây đúng?
	\choice 
		 {Nếu $( P)$ song song với $ a$ thì $( P)$ cũng song song với $ b $}
		{\True Nếu $( P)$ cắt $ a$ thì  $( P)$ cũng cắt $ b $}
		{Nếu $( P)$ chứa $ a$  thì $( P)$ cũng chứa $ b $}
		{Các khẳng định A, B, C đều  sai}
	\loigiai{ 
		Gọi $( Q)\equiv (  a,b)$.
		\begin{itemize}
			\item Khẳng định "Nếu $( P)$ song  song với $ a$ thì $( P)$ cũng song song với $ b $" sai. Vì khi $ b=( P) \cap ( Q) \Rightarrow b \subset ( P)$.
			\item Khẳng định "Nếu $( P)$ chứa $ a$ thì $( P)$ cũng chứa $ b $" sai. Vì khi $( P) \neq ( Q)  \Rightarrow b\parallel ( P)$.
			\item Xét khẳng định "Nếu $( P)$ cắt $ a$  thì $( P)$ cũng cắt $ b $". Giả sử $( P)$ không cắt $ b$ khi đó $ b \subset ( P)$ hoặc $ b\parallel ( P)$. Khi đó, vì $ b \parallel a$ nên $ a \subset ( P)$ hoặc $ a$ cắt $( P)$.
		\end{itemize}
	 }
\end{ex}

% nguồn: Ok_Toan11_TracNghiem_Lan2
\begin{ex}%[1H4N3-1]%[ma:a9de581f]
	Cho hai mặt phẳng $( \alpha)$, $(\beta)$ cắt nhau và cùng  song song với đường thẳng d. Khẳng định nào sau đây là đúng?
	\choice 
		{Giao tuyến của $( \alpha)$, $(  \beta)$ trùng với $ d $}
		{Giao tuyến của $( \alpha)$, $( \beta)$ song song hoặc  trùng với $ d $}
		{Giao tuyến của $( \alpha)$, $( \beta)$ cắt $  d $}
		{\True Giao tuyến của $( \alpha)$, $(  \beta)$ song song với $ d $}
	 \loigiai{
		Do $ d$ không nằm trên mặt phẳng $( \alpha)$, $( \beta)$ nên giao tuyến không thể  trùng với $ d $.\\
		Theo tính chất giao tuyến, ta có giao tuyến song  song với $ d $.
	} 
\end{ex}

% nguồn: Ok_Toan11_TracNghiem_Lan2
\begin{ex}%[1H4N3-1]%[ma:ab03f835]
	Có bao nhiêu mặt phẳng song song với cả hai đường thẳng chéo  nhau?
	\choice 
		 {$1$}
		{$2$} 
		 {$3$}
		 {\True Vô số}
	 \loigiai{
		\begin{center}
			\begin{tikzpicture}[scale=1,line cap=round,line join=round,font=\footnotesize,>=stealth]
					 \def\dai{3.5} \def\rong{2} \def\goc{70}
					 \path (0,0) coordinate (A)++(0:\dai) coordinate (B)
						 ++(\goc:\rong) coordinate (C)++(180:\dai) coordinate (D)
						($(A)+(50:1.2)$) coordinate  (P)++(0:2.7) coordinate (Q)
						($(P)+(90:1.7)$) coordinate  (P')++(0:3.2) coordinate (Q')
						 ($(P)+(-50:0.8)$) coordinate (P1)++(30:2.5) coordinate (Q1);
					 \draw (A)--(B)--(C)--(D)--cycle (P)--node[above,pos=0.9]{$c$}
						 (Q)
						 (P')--node[above,midway]{$a$}
						 (Q')
						(P1)--node[below,pos=0.2]{$b$} 
						(Q1); 
					\draw  pic[draw,angle radius=5mm]{angle=B--A--D};
					 \draw (A) node[shift={(30:0.3)}]{$\alpha$};
				\end{tikzpicture}
		\end{center}
		Gọi $ a$ và $ b$ là $2$ đường  thẳng chéo nhau, $ c$ là đường thẳng song song với $ a$ và cắt $ b$.\\
		Gọi $( \alpha)\equiv ( b,c)$ do $ a \parallel c \Rightarrow a \parallel ( \alpha)$. Giả sử $( \beta) \parallel ( \alpha)$ mà $ b \in ( \alpha) \Rightarrow b \parallel ( \beta)$.\\ 
		Mặt khác: $ a \parallel ( \alpha) \Rightarrow a \parallel ( \beta)$.\\ 
		Có vô số mặt phẳng $( \beta) \parallel ( \alpha)$.  \\
	Vậy có vô số mặt  phẳng song song với $2$ đường thẳng chéo nhau.}
\end{ex}

% nguồn: Ok_Toan11_TracNghiem_Lan2
\begin{ex}%[1H4N3-1]%[ma:7951ed1c]
	Chọn khẳng định  đúng trong các khẳng định sau.
	\choice 
		{Hai đường thẳng phân biệt cùng  song song với một mặt phẳng thì song song với nhau}
		{\True Nếu $a \parallel (P)$ thì tồn  tại trong $(P)$ đường thẳng $b$ để $b \parallel a$}
		{Nếu $\heva{&a \parallel  (P)\\
		 &b\subset (P)}$ thì $a\parallel b$}
		{Nếu $a \parallel (P)$ và đường thẳng $b$ cắt  mặt phẳng $(P)$ thì hai đường thẳng $a$ và $b$ cắt nhau}
	 \loigiai{
	} 
\end{ex}

% nguồn: Ok_Toan11_TracNghiem_Lan2
\begin{ex}%[1H4N3-1]%[ma:acbb49ec]
	Cho hai đường thẳng phân biệt $a$, $b$ và mặt phẳng $\left(\alpha \right)$. Giả sử $a\parallel \left(\alpha \right)$  và $b\parallel \left(\alpha \right)$. Mệnh đề nào sau đây đúng?
	 \choice
		{$a$  và $b$ không có điểm chung}
		{$a$ và  $b$ hoặc song song hoặc chéo nhau}
		 {$a$ và $b$ chéo nhau}
		{\True $a$ và $b$ hoặc song song hoặc  chéo nhau hoặc cắt nhau}
	\loigiai{ 
		 $a$ và $b$ hoặc song song hoặc chéo nhau hoặc cắt nhau.
	 }
\end{ex}

% nguồn: Ok_Toan11_TracNghiem_Lan2
\begin{ex}%[1H4N3-1]%[ma:ff94d356]
	Cho hai đường thẳng phân biệt $a, b$ và mặt  phẳng $\left(\alpha \right)$ thỏa mãn $a\parallel \left(\alpha \right)$ và $b\subset \left(\alpha \right)$. Khi đó
	\choice 
		{$a\parallel b$} 
		{$a$, $b$ chéo nhau} 
		{$a\parallel  b$ hoặc $a$, $b$ chéo nhau}
		{$a$, $b$ cắt  nhau}
	\loigiai{ 
		\begin{center}
			\begin{tikzpicture}[line join=round,line cap=round,scale=0.8]
					 \coordinate (A) at (1,3);
					 \coordinate (B) at (6,3);
					\coordinate (D) at  (0,0);
					 \coordinate (C) at ($(B)+(D)-(A)$);
					\draw (A)--(B)--(C)--(D)--cycle; 
					\draw  (1,4)--(6,4) node[pos=0.2,above]{$a$};
					\draw  (1,1.5)--(5,1.5) node[pos=0.2,above]{$b$};
				\end{tikzpicture}
			\begin{tikzpicture}[line join=round,line cap=round,scale=0.8]
					 \coordinate (A) at (1,3);
					\coordinate (B) at  (6,3);
					 \coordinate (D) at (0,0);
					 \coordinate (C) at ($(B)+(D)-(A)$);
					 \draw (A)--(B)--(C)--(D)--cycle;
					\draw (1,4)--(6,4) node[pos=0.2,above]{$a$}; 
					\draw  (2,1)--(4,2.7) node[pos=0.5,above]{$b$};
					 \draw (1,1.5)--(5,1.5) node[pos=0.2,above]{$c$};
				\end{tikzpicture}
		\end{center}
		Vì $a\parallel (\alpha)$ nên tồn tại đường thẳng $c\subset (\alpha)$ thỏa mãn $a\parallel c$.  Suy ra $b$, $c$ đồng phẳng và xảy ra các trường hợp sau
		\begin{itemize}
			\item Nếu $b$ song song hoặc trùng  với $c$ thì $a\parallel b$.
			\item Nếu $b$ cắt $c$ thì $b$ cắt $(\beta)\equiv (a,c)$ nên $a$, $b$ không đồng phẳng. Do đó $a$, $b$  chéo nhau.
		\end{itemize}
		Vậy  $a\parallel b$ hoặc $a, b$ chéo nhau.
	 }
\end{ex}

% nguồn: Ok_Toan11_TracNghiem_Lan2
\begin{ex}%[1H4N3-1]%[ma:dd0c366d]
	Trong không gian,  vị trí tương đối giữa đường thẳng và mặt phẳng có bao nhiêu trường hợp?
	\choice 
		 {$1$}
		 {$2$}
		{\True $3$} 
		{$4$} 
	 \loigiai{
		Trong không gian, vị trí tương đối giữa đường thẳng và mặt phẳng có $3$ trường hợp đó là đường thẳng cắt mặt phẳng, đường thẳng song song với mặt phẳng và đường thẳng nằm  trong mặt phẳng.
	 }
\end{ex}

% nguồn: Ok_Toan11_TracNghiem_Lan2
\begin{ex}%[1H4N3-1]%[ma:92aba49f]
	Giả thiết nào sau đây  là điều kiện đủ để kết luận đường thẳng $a$ song song với mp $(\alpha)$?
	 \choice
		 {$a \parallel b$ và $b \parallel (\alpha)$}
		{$a \parallel b$ và $b\subset  (\alpha)$}
		{\True $a\cap (\alpha)=\varnothing$} 
		{$a \parallel \left(\beta  \right)$ và $\left(\beta \right) \parallel \left(\alpha \right)$}
	 \loigiai{
		Theo định  nghĩa.
	} 
\end{ex}

% nguồn: Ok_Toan11_TracNghiem_Lan2
\begin{ex}%[1H4N3-1]%[ma:ae640ba8]
	Chọn  khẳng định đúng:
	 \choice
		{Nếu ba mặt phẳng đôi một cắt  nhau theo ba giao tuyến phân biệt thì ba giao tuyến ấy song song với nhau}
		{\True Hai  đường thẳng phân biệt cùng song song với một đường thẳng thứ ba thì song song với nhau}
		{Hai đường thẳng cùng  song song với một đường thẳng thứ ba thì song song với nhau}
		{Nếu ba mặt phẳng đôi một  cắt nhau theo ba giao tuyến thì ba giao tuyến ấy hoặc đồng quy hoặc đôi một song song}
	\loigiai{ 
	Mệnh đề đúng là \lq\lq  Hai đường thẳng phân biệt cùng song song với một đường thẳng thứ ba thì song song với nhau\rq\rq .}
\end{ex}

% nguồn: Ok_Toan11_TracNghiem_Lan2
\begin{ex}%[1H4N3-1]%[ma:29354161]
	Cho đường thẳng $ a$ và mặt phẳng $( P)$ trong không gian. Có bao nhiêu vị  trí tương đối của $ a$ và $( P)$?
	 \choice
		{$2$} 
		{\True  $3$}
		 {$1$}
		{$4$} 
	 \loigiai{
		Có $3$ vị trí tương đối là: $ a$ nằm trong $( P)$, $ a$  song song với $( P)$ và $ a$ cắt $( P)$.
	 }
\end{ex}

% nguồn: Ok_Toan11_TracNghiem_Lan2
\begin{ex}%[1H4N3-1]%[ma:4e08b665]
	Trong không gian, cho hai đường thẳng phân biệt  $ a$, $ b$ và mặt phẳng $( P)$. Những mệnh đề nào sau đây là đúng?
	\choice 
		{\True Nếu $  a$ và $( P)$ có điểm chung thì $ a$ không song song với $( P)$}
		{Nếu $ a$ và  $( P)$ có điểm chung thì $ a$ và $( P)$ cắt nhau}
		{Nếu $ a$ song song với $ b$ và $  b$ nằm trong $( P)$ thì $ a$ song song với $( P)$}
		{Nếu $ a$ và $ b$ song song với  $( P)$ thì $ a$ song song với b}
	\loigiai{ 
		Mệnh đề \lq\lq Nếu $ a$ và $( P)$ có điểm chung thì $ a$  không song song với $( P)$\rq\rq ~ đúng.\\
		Mệnh đề \lq\lq Nếu $ a$ và $( P)$ có điểm chung thì $ a$  và $( P)$ cắt nhau\rq\rq~ sai.\\
		Mệnh đề \lq\lq Nếu $ a$ song song với $ b$  và $ b$ nằm trong $( P)$ thì $ a$ song song với $( P)$\rq\rq ~ sai vì $ a$ và $ b$ có thể cùng thuộc $( P)$.\\
		Mệnh đề \lq\lq Nếu $ a$ và $ b$ song song với $( P)$ thì $ a$ song song với $ b$\rq\rq~  sai.
	 }
\end{ex}

% nguồn: Ok_Toan11_TracNghiem_Lan2
\begin{ex}%[1H4N3-1]%[ma:4435c5e7]
	Cho hai đường thẳng phân biệt $ a$, $b$ và mặt phẳng $( \alpha)$. Giả sử $ a \parallel b$, $ b \parallel ( \alpha)$. Khi  đó
	\choice 
		{$  a \parallel ( \alpha) $}
		{$ a  \subset ( \alpha) $}
		 {$ a$ cắt $( \alpha) $}
		{\True $ a \parallel ( \alpha)$ hoặc $ a  \subset ( \alpha) $}
	 \loigiai{
		Ta có $ a \parallel  ( \alpha)$ hoặc $ a \subset ( \alpha) $.
	 }
\end{ex}

% nguồn: Ok_Toan11_TracNghiem_Lan2
\begin{ex}%[1H4N3-1]%[ma:ce3a6373]
	Hai  đường thẳng cùng song song với đường thẳng thứ ba thì hai đường thẳng đó
	 \choice
		{\True Hoặc song song hoặc trùng nhau} 
		{Chéo  nhau}
		{Trùng  nhau}
		{Song song} 
	\loigiai{ 
		Hai đường thẳng cùng song song với đường thẳng thứ ba thì hai  đường thẳng đó song song hoặc trùng nhau.
	} 
\end{ex}

% nguồn: Ok_Toan11_TracNghiem_Lan2
\begin{ex}%[1H4N3-1]%[ma:1432f443]
	 Cho mặt phẳng $(\alpha)$ và đường thẳng $d\not\subset (\alpha)$. Khẳng định nào sau đây là \textbf{sai}?
	\choice 
		{Nếu $d \parallel (\alpha)$ thì  trong $(\alpha)$ tồn tại đường thẳng $\Delta$ sao cho $\Delta \parallel d$}
		{\True Nếu $d \parallel (\alpha)$ và $b\subset (\alpha)$ thì $b  \parallel d$}
		{Nếu $d\cap (\alpha)=A$ và $d'\subset (\alpha)$  thì $d$ và $d'$ hoặc cắt nhau hoặc chéo nhau}
		 {Nếu $d \parallel c$, $c\subset (\alpha)$ thì $d \parallel (\alpha)$}
	\loigiai{ 
		Nếu $d \parallel (\alpha)$ và $b\subset (\alpha)$ thì $b  \parallel d$ là
		 mệnh đề sai vì $b$ và $d$ có thể chéo nhau.
	 }
\end{ex}

% nguồn: Ok_Toan11_TracNghiem_Lan2
\begin{ex}%[1H4N3-1]%[ma:ab300b6f]
	Khẳng định nào sau đây là đúng? 
	\choice 
		{Nếu đường thẳng $a$ song song với một đường thẳng $b$ nào đó nằm trong $\left(P\right)$ thì $a$  song song với $\left(P\right)$}
		{Nếu đường thẳng $a$ nằm trong mặt phẳng $\left(P\right)$ và song song với một đường thẳng $b$ nào đó nằm trong $\left(P\right)$ thì $a$ song song với $\left(P\right)$} 
		{\True Nếu đường thẳng $a$ không nằm trong mặt phẳng $\left(P\right)$ và song song với một đường thẳng  $b$ nào đó nằm trong $\left(P\right)$ thì $a$ song song với $\left(P\right)$}
		{Nếu đường thẳng $a$ không nằm trong mặt phẳng $\left(P\right)$ thì $a$ song song  với $\left(P\right)$}
	\loigiai{ 
		Ta có mệnh đề đúng là: Nếu đường thẳng $a$ không nằm trong mặt phẳng $\left(P\right)$ và song  song với một đường thẳng $b$ nào đó nằm trong $\left(P\right)$ thì $a$ song song với $\left(P\right)$.
	 }
\end{ex}

% nguồn: Ok_Toan11_TracNghiem_Lan2
\begin{ex}%[1H4N3-1]%[ma:43b90d3c]
	Cho đường thẳng $a$ song song với mặt phẳng $\left(P\right)$ và $b$ là đường thẳng nằm trong $\left(P\right)$. Khi đó trường  hợp nào sau đây không thể xảy ra?
	 \choice
		{$a$ song song $b$} 
		{\True $a$  cắt $b$}
		{$a$ và  $b$ chéo nhau}
		{$a$ và $b$  không có điểm chung}
	\loigiai{ 
		Vì $a\parallel\left(P\right)$ nên $a$ không điểm chung với  mặt phẳng $\left(P\right)$.\\
		Mà $b\subset \left(P\right)$ nên $a$  không điểm chung với $b$ tức $a$ không thể cắt $b$.
	} 
\end{ex}

% nguồn: Ok_Toan11_TracNghiem_Lan2
\begin{ex}%[1H4N3-1]%[ma:a18fe641]
	Mệnh đề nào $sai$ trong các  mệnh đều sau?
	\choice 
		{\True Qua một điểm nằm ngoài một mặt phẳng cho trước có một và chỉ một đường thẳng song song với mặt phẳng  đã cho}
		{Nếu mặt phẳng $\left(\alpha \right)$ chứa hai đường thẳng cắt nhau $a$, $b$ và $a$, $b$ cùng song song với mặt phẳng $\left(\beta \right)$ thì $\left(\alpha \right)$  song song với $\left(\beta \right)$}
		 {Hai mặt phẳng phân biệt cùng song song với mặt phẳng thứ ba thì song song với nhau}
		{Hai mặt phẳng  song song chắn trên hai cát tuyến song song những đoạn thẳng bằng nhau}
	\loigiai{ 
		Qua một điểm nằm ngoài một mặt phẳng cho trước có một  và chỉ một đường thẳng song song với mặt phẳng đã cho.
	 }
\end{ex}

% nguồn: Ok_TracNghiemLop11_Lan1
\begin{ex}%[1H4N3-1]%[ma:c5be7462]
	 Có bao nhiêu mặt phẳng song song với cả hai đường thẳng chéo nhau?
	 \choice
		{1  }
		 {2 }
		{3 } 
		{\True  Vô số}
	\loigiai{ 
		Gọi $a$ và $b$ là 2 đường thẳng chéo nhau, $c$ là  đường thẳng song song với $a$ và cắt $b$.\\
		Gọi $(\alpha) \equiv \left(b,c\right)$.\\ 
		Do $a \parallel c \Rightarrow a  \parallel (\alpha)$.\\
		Giả sử $(\beta) \parallel(\alpha)$. Mà  $b \subset(\alpha) \Rightarrow b \parallel(\beta)$.\\
		Mặt khác, $a \parallel (\alpha) \Rightarrow a \parallel(\beta)$.\\ 
		Có vô số mặt phẳng  $(\beta) \parallel(\alpha)$.\\
		Vậy có  vô số mặt phẳng song song với 2 đường thẳng chéo nhau.
	 }
\end{ex}

% nguồn: Ok_TracNghiemLop11_Lan1
\begin{ex}%[1H4N3-1]%[ma:bfc5c81f]
	Trong không gian,  khẳng định nào dưới đây \textbf{đúng}?
	\choice 
		{Hai đường thẳng phân biệt cùng song song với một mặt phẳng thì song song với nhau} 
		{\True Nếu đường thẳng $a$ không thuộc mặt phẳng $\left( P \right)$ và song song với một đường thẳng thuộc mặt phẳng $\left( P \right)$ thì đường  thẳng $a$ song song với mặt phẳng $\left( P \right)$}
		{Qua một điểm nằm ngoài một mặt phẳng cho trước có một và chỉ một đường thẳng song song với mặt phẳng đã  cho}
		 {Nếu đường thẳng $a$ song song với mặt phẳng $\left( P \right)$ thì đường thẳng $a$ song song với mọi đường thẳng thuộc mặt phẳng $\left( P \right)$}
	\loigiai{ 
		\begin{itemize}
			\item Khẳng định \lq\lq Hai đường thẳng phân biệt cùng song song với một mặt phẳng thì song song với nhau\rq\rq  \,là sai vì hai đường thẳng phân biệt cùng song song với một mặt phẳng có thể cắt nhau hoặc chéo nhau.
			\item Khẳng định \lq\lq Nếu đường thẳng $a$ không thuộc mặt phẳng $\left( P \right)$ và song song với một đường thẳng thuộc mặt phẳng $\left( P \right)$ thì đường thẳng $a$ song song với mặt phẳng $\left( P \right)$\rq\rq \,là sai vì qua một điểm  nằm ngoài một mặt phẳng cho trước có vô số đường thẳng song song với mặt phẳng đã cho.
			\item Khẳng định \lq\lq Nếu đường thẳng $a$ song song với mặt phẳng $\left( P \right)$ thì đường thẳng $a$ song song với  mọi đường thẳng thuộc mặt phẳng $\left( P \right)$\rq\rq\, là sai vì nếu đường thẳng $a$ song song với mặt phẳng $\left( P \right)$ thì đường thẳng $a$ có thể chéo nhau với đường thẳng nào đó thuộc mặt phẳng $\left( P \right)$.
			\item Khẳng đinh \lq\lq Nếu đường thẳng $a$ không thuộc mặt phẳng $\left( P \right)$ và song  song với một đường thẳng thuộc mặt phẳng $\left( P \right)$ thì đường thẳng $a$ song song với mặt phẳng $\left( P \right)$\rq\rq\, là đúng (theo định lý).
		\end{itemize}
	} 
\end{ex}

% nguồn: Ok_TracNghiemLop11_Lan1
\begin{ex}%[1H4N3-1]%[ma:f1fa43db]
	Trong không gian, hai đường  thẳng được gọi là song song với nhau nếu?
	\choice 
		{chúng  không có điểm chung}
		{\True chúng cùng thuộc một mặt phẳng và không có  điểm chung}
		{chúng  có 1 điểm chung}
		{không cùng thuộc một  mặt phẳng}
	\loigiai{ 
		Trong không gian, hai đường  thẳng được gọi là song song với nhau nếu chúng cùng thuộc một mặt phẳng và không có điểm chung.
	 }
\end{ex}

% nguồn: Ok_TracNghiemLop11_Lan1
\begin{ex}%[1H4N3-1]%[ma:629cdd75]
	Với điều kiện nào sau đây thì đường thẳng $a$  song song với mặt phẳng $(\alpha)$.
	\choice 
		{$a\parallel b$ và $b\cap  (\alpha)=\varnothing$}
		{$a\parallel b$  và $b\parallel (\alpha)$}
		 {\True $a\parallel b$ và $b\subset (\alpha)$}
		{$a\cap  (\alpha)=\varnothing$}
	 \loigiai{Lý thuyết}
\end{ex}

% nguồn: Ok_Toan11_TracNghiem_Lan2
\begin{ex}%[1H4H3-1]%[ma:ac6290bc]
	Cho các giả thiết sau đây. Giả thiết nào kết luận đường thẳng $a$  song song với mặt phẳng $\left(\alpha \right)$?
	\choice 
		{$a\parallel b$ và $b\subset \left(\alpha  \right)$}
		{$a\parallel  \left(\beta \right)$ và $\left(\beta \right)\parallel \left(\alpha \right)$}
		{$a\parallel b$ và  $b\parallel \left(\alpha \right)$}
		{\True $a\cap \left(\alpha  \right)=\varnothing$}
	 \loigiai{
		Chọn $a\cap \left(\alpha \right)=\varnothing$. 
	} 
\end{ex}

% nguồn: Ok_Toan11_TracNghiem_Lan2
\begin{ex}%[1H4H3-1]%[ma:541c138c]
	Cho ba đường thẳng đôi một chéo nhau $a$, $b$, $c$. Gọi $\left(P\right)$ là mặt phẳng qua $a$, $\left(Q\right)$ là mặt phẳng qua $b$ sao cho giao tuyến của $\left(P\right)$ và $\left(Q\right)$ song song với $c$. Có nhiều nhất bao nhiêu mặt phẳng $\left(P\right)$ và $\left(Q\right)$  thỏa mãn yêu cầu trên?
	 \choice
		{Vô số mặt  phẳng $\left(P\right)$ và $\left(Q\right)$}
		{Một mặt phẳng  $\left(P\right)$, vô số mặt phẳng $\left(Q\right)$}
		 {Một mặt phẳng $\left(Q\right)$, vô số mặt phẳng $\left(P\right)$}
		{\True  Một mặt phẳng $\left(P\right)$, một mặt phẳng $\left(Q\right)$}
	\loigiai{ 
		\begin{center}
			\begin{tikzpicture}[scale=1,font=\footnotesize,line join=round,line cap=round,>=stealth,declare function={a=3; b=5; goc=40; }]
					 \path
						(0,0) coordinate (A) 
						++(goc:a) coordinate (B) 
						++(90:b) coordinate  (C)
						($(C)-(B)$) coordinate  (D)
						(A)++(1.5*a,0) coordinate (E) 
						($(C)+(E)-(B)$) coordinate (F); 
					\draw  (A)--(B)--(C)--(D)--cycle (C)--(F)--(E)--(B);
					\draw  ($(B)+(-0.6,0.8)$)-- ($(D)+(0.4,-0.3)$) node[right]{$a$}
						($(B)+(0.6,0.5)$)-- ($(F)+(-0.5,-0.3)$)  node[left]{$b$}
						($(E)+(1,0.5)$)-- ($(F)+(1,-0.5)$) node[right]{$c$}; 
					\foreach  \x/\y/\z in {B/A/D}{
						\path  (\y) pic[draw,angle radius=35pt,angle eccentricity=0.7,"$(P)$"]{angle = \x--\y--\z};
					 }
					\foreach \x/\y/\z in {F/E/B}{ 
						\path  (\y) pic[draw,angle radius=35pt,angle eccentricity=0.7,"$(Q)$"]{angle = \x--\y--\z};
					 }
				\end{tikzpicture}
		\end{center}
		Vì $c$ song song với giao tuyến  của $\left(P\right)$ và $\left(Q\right)$ nên $c\parallel \left(P\right)$ và $c\parallel \left(Q\right)$.\\
		Khi đó, $\left(P\right)$ là mặt phẳng chứa $a$ và song song với $c$, mà $a$ và $c$ chéo nhau nên chỉ có một mặt  phẳng như vậy.\\
		Tương  tự cũng chỉ có một mặt phẳng $\left(Q\right)$ chứa $b$ và song song với $c$.\\
		Vậy có nhiều nhất một mặt phẳng $\left(P\right)$ và một mặt phẳng $\left(Q\right)$ thỏa  yêu cầu bài toán.
	 }
\end{ex}

% nguồn: Ok_Toan11_TracNghiem_Lan2
\begin{ex}%[1H4H3-1]%[ma:67cc61b8]
	Tìm khẳng định \textbf{sai} trong các khẳng định  sau đây
	 \choice
		{Nếu hai mặt phẳng song song cùng cắt mặt phẳng thứ ba thì hai giao tuyến  tạo thành song song với nhau}
		{Ba mặt phẳng đôi một song song chắn trên hai đường thẳng  chéo nhau những đoạn thẳng tương ứng tỉ lệ}
		 {Nếu mặt phẳng $\left(P\right)$ song song với mặt phẳng $\left(Q\right)$ thì mọi đường thẳng nằm trên mặt phẳng $\left(P\right)$ đều song song với mặt phẳng $\left(Q\right)$}
		{\True Nếu mặt  phẳng $\left(P\right)$ có chứa hai đường thẳng phân biệt và hai đường thẳng đó cùng song song song với mặt phẳng $\left(Q\right)$ thì mặt phẳng $\left(P\right)$ song song với mặt phẳng $\left(Q\right)$}
	 \loigiai{
		\begin{center}
			\begin{tikzpicture}[scale=1,font=\footnotesize,line join=round,line cap=round,>=stealth,declare function={gocx=80; goc=-150; a=5; b=a/2; h=4;}]
					 \path (0,0) coordinate (A)--+(gocx:h) coordinate (S)
						(a,0) coordinate (D) 
						(goc:b) coordinate  (B)
						+(a,0) coordinate  (C);
					\path ($(A)!0.5!(S)$) coordinate (M); 
					 \path ($(D)!0.5!(S)$) coordinate (N);
					 \path ($(A)!1/3!(S)$) coordinate (P);
					\path ($(D)!1/3!(S)$)  coordinate (Q);
					\draw[dashed] (S)--(A) 
						 (D)--(A)--(B)
						 (M)--(N)
						 (P)--(Q);
					\draw  (S)--(D)--(C)--(S)--(B)--(C)--cycle;
					 \foreach \x/\goc in {S/90,A/150,B/180,C/-60,D/0,M/180,P/180,N/30,Q/30}
					\draw[fill=white] (\x) node[shift={(\goc:7pt)}]{$\x$} circle  (1pt);
				\end{tikzpicture}
		\end{center}
		 Ví dụ $\left(SAD\right)$ chứa $MN$; $PQ$ cùng song song với $\left(ABCD\right)$ nhưng $\left(SAD\right)$ cắt $\left(ABCD\right)$.
	 }
\end{ex}

% nguồn: Ok_Toan11_TracNghiem_Lan2
\begin{ex}%[1H4H3-1]%[ma:9d794c40]
	Trong các khẳng định  sau khẳng định nào \textbf{sai}?
	 \choice
		 {\True Nếu một đường thẳng song song với một trong hai mặt phẳng song song thì nó song song với mặt phẳng còn lại}
		{Nếu một đường thẳng cắt một trong hai mặt phẳng song song thì nó cắt  mặt phẳng còn lại}
		{Nếu hai đường thẳng song song thì chúng cùng nằm trên một mặt  phẳng}
		{Nếu hai mặt phẳng phân biệt cùng song song với một mặt phẳng thì chúng  song song với nhau}
	\loigiai{ 
		Giả  sử $\left(\alpha \right)$ song song với $\left(\beta \right)$. Một đường thẳng $a$ song song với $\left(\beta \right)$ có thể nằm trên $\left(\alpha \right)$.
	} 
\end{ex}

% nguồn: Ok_Toan11_TracNghiem_Lan2
\begin{ex}%[1H4H3-1]%[ma:aa24f3b5]
	Cho hai mặt phẳng $\left(P\right)$, $\left(Q\right)$ cắt nhau theo giao tuyến là đường thẳng $d$. Đường thẳng $a$ song song với cả hai mặt phẳng $\left(P\right)$, $\left(Q\right)$. Khẳng định nào  sau đây đúng?
	\choice 
		{$a$, $d$ trùng  nhau}
		{$a$, $d$  chéo nhau}
		 {\True $a$ song song $d$}
		{$a$,  $d$ cắt nhau}
	\loigiai{ 
		Sử dụng hệ quả: Nếu hai mặt phẳng phân biệt cùng song song với  một đường thẳng thì giao tuyến của chúng cũng song song với đường thẳng đó.
	 }
\end{ex}

% nguồn: Ok_Toan11_TracNghiem_Lan2
\begin{ex}%[1H4H3-1]%[ma:49a2e2fb]
	Cho lăng trụ $ABC.A'B'C'$. Gọi $G$,  $G'$ lần lượt là trọng tâm các tam giác $ABC$ và $A'B'C'$, $M$ là điểm trên cạnh $AC$ sao cho $AM=2MC$. Khẳng định nào sau đây \textbf{sai}?
	 \choice
		{$GG'  \parallel \left(ABB'A'\right)$}
		{$GG' \parallel  \left(ACC'A'\right)$}
		{\True Đường thẳng $MG'$  cắt mặt phẳng $\left(BCC'B'\right)$}
		 {$\left(MGG'\right) \parallel \left(BCC'B'\right)$}
	 \loigiai{
		\begin{center}
			\begin{tikzpicture}[scale=0.8,font=\footnotesize,line join=round,line cap=round,>=stealth]
					\coordinate (A) at  (-2,0);
					\coordinate  (B) at (-1,-2);
					\coordinate (C) at (2,0); 
					\coordinate (H) at  (1,4);
					\foreach \i in {A,B,C}{\coordinate (\i') at ($(\i)+(H)$); \draw  (\i)--(\i');}
					\draw (A')--(B')--(C')--cycle; 
					\draw (A)--(B)--(C); 
					\draw[dashed,thin](A)--(C); 
					\coordinate (N)  at ($(B)!0.5!(C)$);
					\coordinate (N') at  ($(B')!0.5!(C')$);
					\coordinate (G)  at ($(A)!0.67!(N)$);
					 \coordinate (G') at ($(A')!0.67!(N')$);
					 \coordinate (M) at ($(A)!0.67!(C)$);
					\draw  (A')--(N')--(N);
					 \draw[dashed,thin](A)--(N)
						(G')--(G)--(M); 
					\foreach \i/\g in {A/180,B/-90,C/0,A'/180,B'/180,C'/0,M/90,N/-10,N'/-10,G/-90,G'/90}{\draw[fill=black](\i)  circle (1pt)
						($(\i)+(\g:4mm)$)  node[scale=1]{$\i$};}
				\end{tikzpicture}
		\end{center}
		Ta có $GG' \parallel AA'$ và $MG \parallel BC$ nên $GG'  \parallel \left(ACC'A'\right)$ là mệnh đề đúng.\\
		$GG'  \parallel \left(ABB'A'\right)$ là mệnh đề đúng.\\
		$\left(MGG'\right) \parallel  \left(BCC'B'\right)$ là mệnh đề đúng.\\
		 Đường thẳng $MG'$ cắt mặt phẳng $\left(BCC'B'\right)$ là mệnh đề sai.
	} 
\end{ex}

% nguồn: Ok_Toan11_TracNghiem_Lan2
\begin{ex}%[1H4H3-1]%[ma:7f0b74b5]
	Trong các khẳng định sau,  khẳng định nào là \textbf{đúng}?
	\choice 
		{Nếu đường thẳng $d$ song song mặt phẳng $(P)$ thì trong $(P)$ có duy nhất một đường thẳng $a$ song song với  $d$}
		{Nếu đường thẳng $d$ song song với  mặt phẳng $(P)$ thì $d$ song song với mọi đường thẳng nằm trong $(P)$}
		{\True Nếu đường thẳng $d$ song song với mặt phẳng $(P)$ thì trong $(P)$ tồn tại  đường thẳng $a$ song song với $d$}
		{Nếu đường thẳng $d$ song song mặt phẳng $(P)$, đường thẳng  $a$ bất kỳ nằm trong $(P)$ thì $a$ và $d$ chéo nhau}
	\loigiai{ 
		Khẳng định đúng là \lq\lq Nếu đường thẳng $d$ song song với mặt phẳng $(P)$  thì trong $(P)$ tồn tại đường thẳng $a$ song song với $d$\rq\rq.
	 }
\end{ex}

% nguồn: Ok_Toan11_TracNghiem_Lan2
\begin{ex}%[1H4H3-1]%[ma:3f0da6c4]
	Cho các mệnh đề  sau:
	\begin{itemize}
		\item Nếu $a \parallel \left(P\right)$ thì $a$ song song với  mọi đường thẳng nằm trong $\left(P\right)$.
		\item  Nếu $a \parallel \left(P\right)$ thì $a$ song song với một đường thẳng nào đó nằm trong $\left(P\right)$.
		 \item Nếu $a\parallel \left(P\right)$ thì có vô số đường thẳng nằm trong $\left(P\right)$ song song với $a$.
		\item Nếu $a\parallel \left(P\right)$ thì có  một đường thẳng $d$ nào đó nằm trong $\left(P\right)$ sao cho $a$ và $d$ đồng phẳng.
	\end{itemize}
	Số mệnh đề đúng là 
	 \choice
		 {$2$}
		{\True  $3$}
		{$4$} 
		{$1$} 
	 \loigiai{
		\begin{itemize}
			\item Nếu $a\parallel \left(P\right)$ thì $a$ song song với mọi đường thẳng nằm  trong $\left(P\right)$ (Sai).
			\item Nếu $a\parallel \left(P\right)$ thì $a$ song song với một đường  thẳng nào đó nằm trong $\left(P\right)$ (Đúng).
			 \item Nếu $a\parallel \left(P\right)$ thì có vô số đường thẳng nằm trong $\left(P\right)$ song song với $a$ (Đúng).
			\item Nếu $a\parallel \left(P\right)$ thì có một  đường thẳng $d$ nào đó nằm trong $\left(P\right)$ sao cho $a$ và $d$ đồng phẳng (Đúng).
		\end{itemize}
		 Vậy có 3 mệnh đề đúng.
	} 
\end{ex}

% nguồn: Ok_Toan11_TracNghiem_Lan2
\begin{ex}%[1H4H3-1]%[ma:61d231b6]
	 Trong không gian, đường thẳng $a$ song song với mặt phẳng $(P)$ nếu
	 \choice
		{$a \not \subset(P)$} 
		{$\heva{&a \parallel b \\ &b  \subset(P)}$}
		{$\heva{&a  \parallel b \\ &b \not \subset(P)}$}
		 {\True $\heva{&a \parallel b \\ &b \subset(P) \\ &a \not \subset(P)}$}
	\loigiai{ 
		Đường thẳng $a$ song song với mặt phẳng $(P)$ khi và chỉ khi $a$ không nằm trong $(P)$, đồng thời $a$ song song với một đường thẳng $b$ nằm trong  $(P)$.
	} 
\end{ex}

% nguồn: Ok_Toan11_TracNghiem_Lan2
\begin{ex}%[1H4H3-1]%[ma:d1cdc87a]
	Cho ba đường thẳng đôi một chéo nhau $ a$, $b$, $c$. Gọi $( P)$ là mặt phẳng qua $ a$, $( Q)$ là mặt phẳng qua $ b$ sao cho giao tuyến của $( P)$ và $( Q)$ song song với $ c$. Có nhiều nhất bao nhiêu mặt phẳng $( P)$ và $( Q)$ thỏa mãn  yêu cầu trên?
	 \choice
		{\True Một mặt phẳng $(  P)$, một mặt phẳng $( Q) $}
		 {Một mặt phẳng $( P)$, vô số mặt phẳng $( Q) $}
		{Một mặt phẳng $( Q)$, vô số mặt phẳng $( P)  $}
		{Vô  số mặt phẳng $( P)$ và $( Q) $}
	 \loigiai{
		\begin{center}
			\begin{tikzpicture}[scale=1,line cap=round,line join=round,font=\footnotesize,>=stealth]
					\def\cao{4.5} \def\ngang{2.5} \def\goc{30} 
					\path (0,0) coordinate (A)++(\goc:\ngang)  coordinate (B)
						++(90:\cao)  coordinate (C)++(-180+\goc:\ngang) coordinate (D)
						(B)++(-\goc:\ngang) coordinate (E)++(90:\cao) coordinate (F) 
						 ($(B)+(140:0.5)$) coordinate (P)++(120:2.7) coordinate (Q)
						($(B)+(20:0.5)$) coordinate (P')++(70:2.7) coordinate  (Q')
						($(B)+(-20:3.5)$) coordinate (P1)++(90:4.2) coordinate (Q1); 
					\draw (A)--(B)--(C)--(D)--cycle (B)--(E)--(F)--(C) 
						(P)--node[left,pos=0.5]{$a$} 
						(Q) 
						 (P')--node[right,pos=0.8]{$b$}
						 (Q')
						(P1)--node[right,pos=0.8]{$c$} 
						(Q1); 
					\draw  pic[draw,angle radius=7mm]{angle=B--A--D};
					\draw (A)  node[shift={(55:0.4)}]{$P$};
					\draw  pic[draw,angle radius=7mm]{angle=F--E--B};
					\draw  (E) node[shift={(120:0.4)}]{$Q$};
				\end{tikzpicture}
		\end{center}
		Vì $ c$ song song với giao tuyến của $( P)$ và $( Q)$ nên $ c\parallel ( P)$ và  $ c\parallel ( Q)$.\\
		Khi đó, $( P)$ là mặt phẳng chứa $ a$ và song song với $ c$, mà $ a$ và $ c$ chéo nhau  nên chỉ có một mặt phẳng như vậy.\\
		Tương tự cũng chỉ có một mặt phẳng $( Q)$ chứa $ b$ và song  song với $ c$.\\
		Vậy có nhiều nhất một mặt phẳng  $( P)$ và một mặt phẳng $( Q)$ thỏa yêu cầu bài toán.
	} 
\end{ex}

% nguồn: Ok_Toan11_TracNghiem_Lan2
\begin{ex}%[1H4H3-1]%[ma:1f6258ff]
	Trong các mệnh  đề sau, mệnh đề nào đúng?
	 \choice
		{Hai  đường thẳng cùng song song với một mặt phẳng thì song song với nhau}
		{Hai đường thẳng cùng song song với một mặt phẳng thì trùng  nhau}
		{Hai đường thẳng cùng song song với một mặt phẳng thì chéo  nhau}
		{\True Hai đường thẳng cùng song song với một mặt phẳng có thể chéo nhau, song song, cắt nhau hoặc trùng nhau} 
	 \loigiai{
		Lý thuyết: Hai đường  thẳng cùng song song với một mặt phẳng có thể chéo nhau, song song, cắt nhau hoặc trùng nhau.
	 }
\end{ex}

% nguồn: Ok_TracNghiemLop11_Lan1
\begin{ex}%[1H4H3-1]%[ma:71c63152]
	Cho mặt phẳng $\left( P \right)$ và điểm $A$ không thuộc mặt phẳng $\left( P \right)$. Số đường thẳng qua $A$ và  song song với mặt phẳng $\left( P \right)$ là
	 \choice
		{$0$} 
		{\True Vô số} 
		 {$1$}
		{$2$} 
	 \loigiai{
		Cho mặt phẳng $\left( P \right)$ và điểm $A$ không thuộc mặt phẳng $\left( P \right)$. Số  đường thẳng qua $A$ và song song với mặt phẳng $\left( P \right)$ là ô số.
	} 
\end{ex}

% nguồn: Ok_TracNghiemLop11_Lan1
\begin{ex}%[1H4H3-1]%[ma:d386eeac]
	Tìm khẳng định đúng trong  các khẳng định sau.
	\choice 
		{\True Nếu một đường thẳng song song với một  mặt phẳng thì nó song song với một đường thẳng nào đó nằm trong mặt phẳng đó}
		{Nếu hai mặt  phẳng cùng song song với mặt phẳng thứ ba thì chúng song song với nhau}
		{Nếu ba mặt phẳng phân biệt đôi  một cắt nhau theo ba giao tuyến thì ba giao tuyến đó phải đồng quy}
		{Trong không  gian, hai đường thẳng cùng vuông góc với đường thẳng thứ ba thì hai đường thẳng đó song song với nhau}
	\loigiai{ 
		\begin{itemize}
			 \item Nếu một đường thẳng song song với một mặt phẳng thì nó song song với một đường thẳng nào đó nằm trong mặt phẳng đó là mệnh đề đúng.
			\item Nếu hai  mặt phẳng cùng song song với mặt phẳng thứ ba thì chúng song song với nhau sai vì hai mặt phẳng có thể trùng nhau.
			\item  Nếu ba mặt phẳng phân biệt đôi một cắt nhau theo ba giao tuyến thì ba giao tuyến đó phải đồng quy sai vì ba giao tuyến có thể song song hoặc trùng nhau.
			\item  Trong không gian, hai đường thẳng cùng vuông góc với đường thẳng thứ ba thì hai đường thẳng đó song song với nhau sai vì hai đường thẳng đó có thể trùng nhau hoặc chéo nhau hoặc cắt nhau.
		\end{itemize}
	} 
\end{ex}

% nguồn: Ok_TracNghiemLop11_Lan1
\begin{ex}%[1H4H3-1]%[ma:2648dd9d]
	Cho hai mặt phẳng $(P)$ và $(Q)$ cắt nhau theo giao tuyến là đường thẳng $a$. Đường thẳng $d$ song song với cả hai mặt phẳng $(P)$  và $(Q)$. Khẳng định nào sau đây đúng?
	\choice 
		{$a$, $d$ trùng  nhau}
		 {$a$, $d$ chéo nhau}
		{\True  $a$, $d$ song song}
		{$a$,  $d$ cắt nhau}
	\loigiai{ 
		Sử dụng hệ quả:  Nếu hai mặt phẳng phân biệt cùng song song với một đường thẳng thì giao tuyến của chúng (nếu có) cũng song song với đường thẳng đó.
	 }
\end{ex}


%%%%%%%%%%%%%%%% ĐÚNG/SAI %%%%%%%%%%%%%%%%

% nguồn: Ok_dungsai_lop11_Lan1
\begin{ex}%[1H4N3-1]%[ma:67d957c9]
 Cho mặt phẳng $\left(  \alpha \right)$ và đường thẳng $d\not\subset\left( \alpha \right)$. Khi đó
  \choiceTF
  {Nếu $d\parallel \left(  \alpha \right)$ và $a\subset \left( \alpha \right)$ thì $a\parallel d$}
  {\True Nếu  $d\parallel a$ và $a\subset \left( \alpha \right)$ thì $d\parallel \left( \alpha \right)$}
  {\True Nếu $d\parallel \left( \alpha \right)$ thì trong $\left( \alpha \right)$ tồn tại đường thẳng $a$ sao cho  $a\parallel d$ }
  {Nếu  $d\parallel \left( \alpha \right)$ và $a\parallel \left( \alpha \right)$ thì $a\parallel d$ }
  \loigiai{
  \begin{itemchoice}
   \itemch \textbf{Sai}.\\ 
   Nếu $d\parallel \left( \alpha \right)$  và $a\subset \left( \alpha \right)$ thì $a\parallel d$ hoặc $a$ chéo $d$.
   \itemch \textbf{Đúng}.\\ 
    Theo dấu hiệu nhận biết đường thẳng song song mặt phẳng.
   \itemch \textbf{Đúng}.\\ 
    Theo dấu hiệu nhận biết đường thẳng song song mặt phẳng.
   \itemch \textbf{Sai}.\\ 
   Vì $a$ và $d$ có thể trùng nhau, song song hoặc cắt nhau. 
  \end{itemchoice}
  }
\end{ex}

% nguồn: Ok_dungsai_lop11_Lan1
\begin{ex}%[1H4N3-1]%[ma:ea989ee5]
 Cho hai mặt phẳng $(P)$ và $(Q)$ phân biệt, hai đường thẳng  song song $a$ và $b$. Khi đó
  \choiceTF
  {\True Chỉ có một mặt phẳng chứa  đường thẳng $a$ và đường thẳng $b$}
  {\True Nếu mặt phẳng $(P)$ chứa đường thẳng $a$ và mặt phẳng $(Q)$ chứa  đường thẳng $b$ thì giao tuyến của $(P)$ và $(Q)$ song song với đường thẳng $a$}
  {\True Nếu $a\subset(P)$ và  $b\subset (Q)$ thì $a\parallel(Q)$}
  {Nếu $(P)$ chứa đường thẳng  $a$ thì $(P)$ cũng chứa đường thẳng $b$}
  \loigiai{
  \begin{itemchoice}
   \itemch \textbf{Đúng}.\\ Có một và chỉ một mặt phẳng chứa đường thẳng $a$ và  đường thẳng $b$.
   \itemch \textbf{Đúng}.\\ Vì $a\parallel b$ nên giao tuyến của  $(P)$ và $(Q)$ song song với đường thẳng $a$
   \itemch \textbf{Đúng}.\\ $\heva{& a\subset(P)\\&b\subset (Q)\\&a\parallel b}\Rightarrow a\parallel(Q)$. 
   \itemch \textbf{Sai}.\\ Có vô số mặt phẳng chứa đường thẳng $a$ mà không chứa đường thẳng $b$  ($a$, $b$ là hai đường thẳng song song).
 \end{itemchoice}}
\end{ex}

% nguồn: Ok_dungsai_lop11_Lan1
\begin{ex}%[1H4N3-1]%[ma:166e76af]
 Cho mặt phẳng $(P)$ và hai đường thẳng  song song $a$ và $b$. Khi đó
 \choiceTF 
   {\True Có vô số mặt phẳng chứa đường thẳng $a$ mà không chứa đường thẳng $b$}
  {Nếu mặt phẳng $(P)$ song song với đường thẳng $a$ thì mặt phẳng $(P)$ cũng song song với đường thẳng  $b$}
  {\True Nếu mặt phẳng $(P)$ cắt đường thẳng $a$  thì mặt phẳng $(P)$ cũng cắt đường thẳng $b$}
  {Nếu mặt phẳng  $(P)$ chứa đường thẳng $a$ thì mặt phẳng $(P)$ cũng chứa đường thẳng $b$}
 \loigiai{ 
  \begin{itemchoice}
   \itemch \textbf{Đúng}. Gọi $M$ là điểm tùy ý không thuộc  $b$. Khi đó có vô số mặt phẳng $(M, a)$ chứa $a$ và không chứa đường thẳng $b$.
   \itemch \textbf{Sai}. Mặt phẳng $(P)$ có thể chứa đường thẳng $b$. 
   \itemch \textbf{Đúng}. Dựa vào tính chất:  Hai đường thẳng song song, mặt phẳng cắt đường thẳng này thì cũng cắt đường thẳng kia.
   \itemch \textbf{Sai}. Mặt phẳng $(P)$ có thể song song với  đường thẳng $b$.
  \end{itemchoice}
 } 
\end{ex}

% nguồn: Ok_dungsai_lop11_Lan1
\begin{ex}%[1H4H3-1]%[ma:74185daa]
 Cho hai đường  thẳng $a$, $b$ và mặt phẳng $(\alpha)$ bất kì.
 \choiceTF 
  {\True Nếu  $a \subset(\alpha)$; $b \not \subset (\alpha)$; $b \parallel a$ thì $b \parallel (\alpha)$}
  {\True Nếu $a \subset (P)$ và $b  \parallel(P)$ thì $a, b$ không có điểm chung}
  {Nếu $a \subset (P)$ và $b \parallel(P)$ thì $a, b$ chéo  nhau}
  {Nếu $a \subset(\alpha)$; $b \not  \subset (\alpha)$; $b$ cắt $(\alpha)$ thì $b$ cắt $a$}
 \loigiai{ 
  Dựa  vào tính chất ta có
  \begin{itemchoice}
   \itemch \textbf{Đúng.} Vì Nếu $a \subset  (P)$ và $b \parallel(P)$ thì $a, b$ không có điểm chung.
   \itemch \textbf{Đúng.} Vì Nếu $a \subset(\alpha)$; $b \not \subset (\alpha)$; $b \parallel a$ thì  $b \parallel (\alpha)$.
   \itemch  \textbf{Sai.} Vì Nếu $a \subset(\alpha)$; $b \not \subset (\alpha)$; $b$ cắt $(\alpha)$ thì $b$ cắt $a$.
   \itemch \textbf{Sai.} Vì Nếu  $a \subset (P)$ và $b \parallel(P)$ thì $a, b$ chéo nhau.
  \end{itemchoice}
 } 
\end{ex}

% nguồn: Ok_dungsai_lop11_Lan2
\begin{ex}%[1H4H3-1]%[ma:d184d120]
	Cho mặt phẳng $(P)$ và hai  đường thẳng song song $a$ và $b$
	 \choiceTF
		{\True Có vô số mặt phẳng  chứa đường thẳng $a$ mà không chứa đường thẳng $b$}
		{Nếu mặt phẳng $(P)$ song song với  đường thẳng $a$ thì mặt phẳng $(P)$ cũng song song với đường thẳng $b$}
		{\True Nếu mặt phẳng $(P)$ cắt đường thẳng $a$ thì mặt  phẳng $(P)$ cũng cắt đường thẳng $b$}
		{Nếu mặt phẳng  $(P)$ chứa đường thẳng $a$ thì mặt phẳng $(P)$ cũng chứa đường thẳng $b$}
	 \loigiai{
		\begin{itemchoice}
			\itemch Có vô số mặt phẳng chứa đường thẳng $a$ mà  không chứa đường thẳng $b$.
			\itemch Nếu mặt phẳng $(P)$ song song với đường thẳng $a$ thì mặt phẳng $(P)$ có thể song song hoặc chứa  đường thẳng $b$.
			\itemch Nếu mặt phẳng $(P)$ cắt đường thẳng $a$ thì mặt phẳng $(P)$ cũng  cắt đường thẳng $b$.
			\itemch Có vô số mặt phẳng  chứa đường thẳng $a$ mà không chứa đường thẳng $b$ ($a$, $b$ là hai đường thẳng song song).
		\end{itemchoice}
	 }
\end{ex}


%%%%%%%%%%%%%%%% TRẢ LỜI NGẮN %%%%%%%%%%%%%%%%

% nguồn: Ok_TraLoiNgan_Lop11_Lan1
\begin{ex}%[1H4V3-1]%[ma:5155db46]
	Cho tứ diện $ABCD$. Gọi $M$ là một điểm bất kì trên cạnh $BC$ sao cho $BM=3MC$, $\left(\alpha\right)$ là mặt phẳng qua $M$ và song song với $AB$ và $CD$, cắt các cạnh $BD$,  $AD$, $AC$ lần lượt tại $N$, $P$, $Q$. Tính tỷ số $\dfrac{MN}{PQ}$.
	 \shortans{1}
	 \loigiai{
	\immini{  Ta có $
		\left.\begin{array}{l}
		A B \parallel (\alpha) \\ 
		(A B  C) \supset A B \\
		(A B C) \cap(\alpha)=M Q 
	\end{array}\right\}\Rightarrow M Q \parallel AB\quad(1)
	 $.
	Tương tự,  ta có $NP \parallel AB\quad(2)$.\\
	Lại có		$\left.\begin{array}{l}
	 C D \parallel (\alpha) \\
	(A C D) \supset  C D \\
	(A C D)  \cap(\alpha)=PQ
\end{array}\right\}\Rightarrow PQ \parallel CD\quad(3)$
Tương tự, ta có $M N \parallel C  D\quad(4)$\\
Từ (1) và (2) suy ra $M  Q \parallel N P\quad(5)$.\\
Từ (3) và (4) suy  ra $P Q \parallel M N\quad(6)$.\\
Từ (5) và (6) suy ra  $MNPQ$ là hình bình hành.\\
Suy ra $\dfrac{MN}{PQ}=1$.}{\begin{tikzpicture}[line join = round,line cap = round,font=\footnotesize,>=stealth,scale=1]
		\path 
			(0:0) coordinate (B) 
			(70:3.5) coordinate (A) 
			(0:5) coordinate  (C)
			(-20:3.5) coordinate (D) 
			($(B)!.6667!(C)$) coordinate (M) 
			($(B)!.6667!(D)$) coordinate  (N)
			($(D)!.3333!(A)$) coordinate (P) 
			 ($(A)!.6667!(C)$) coordinate (Q);
		\draw  (A)--(B)--(D)--(C)--cycle (A)--(D)
			(N)--(P)--(Q); 
		\draw[dashed] (B)--(C) 
			(N)--(M)--(Q); 
		\foreach  \x/\g in {A/90,B/180,D/-90,C/30,M/-50,N/-100,P/150,Q/66}\draw[fill=black] (\x) circle (.05)+(\g:.3) node{$\x$};
	\end{tikzpicture}
 }
 }
\end{ex}


%%%%%%%%%%%%%%%% TỰ LUẬN %%%%%%%%%%%%%%%%

% nguồn: Ok_TuanLuan_Toan11_Lan2
\begin{ex}%[1H4N3-1]%[ma:240cf1b1]
	Cho ba đường thẳng $a$, $b$, $c$ đôi một song song với nhau và không cùng nằm trong một mặt phẳng. Chứng minh rằng đường thẳng $ a$ song song với mặt phẳng $( b,c)$. 
	 \loigiai{
		\begin{center}
			\begin{tikzpicture}[scale=1,line cap=round,line join=round,font=\footnotesize,>=stealth]
					\def\dai{3.5} \def\rong{2}  \def\goc{70}
					\path  (0,0) coordinate (A)++(0:\dai) coordinate (B)
						 ++(\goc:\rong) coordinate (C)++(180:\dai) coordinate (D)
						($(A)+(50:1.5)$) coordinate (P)++(5:2.7) coordinate (Q) 
						 ($(P)+(90:1.5)$) coordinate (P')++(5:2.7) coordinate (Q')
						($(P)+(-100:0.8)$) coordinate (P1)++(5:2.7) coordinate (Q1); 
					\draw (A)--(B)--(C)--(D)--cycle (P)--node[above,midway]{$b$} 
						(Q) 
						(P')--node[above,midway]{$a$} 
						 (Q')
						(P1)--node[above,midway]{$c$} 
						(Q1); 
				\end{tikzpicture}
		\end{center}
		Ba đường thẳng $a$, $b$, $c$ không cùng nằm trong một mặt phẳng nên đường thẳng $ a$ không nằm trong $ mp( b,c)$. Vì đường thẳng $ a$ song song với đường thẳng $ b$ và đường thẳng $ b$ nằm trong $ mp( b,c)$ nên đường thẳng $  a$ song song với mặt phẳng $( b,c)$.
	} 
\end{ex}

% nguồn: Ok_TuanLuan_Toan11_Lan2
\begin{ex}%[1H4H3-1]%[ma:541b6b63]
	Cho hai điểm $A$, $B$ cùng thuộc mặt phẳng $( P)$ và một điểm $C$ không thuộc $( P)$. Vẽ đường thẳng ${d_1}$ đi qua $A$, $B$; ${d_2}$ đi qua $C$ và song song $AB$. Tìm số điểm  chung của mỗi đường thẳng vừa vẽ với $( P)$. Xét vị trí tương đối của mặt phẳng $( P)$ lần lượt đối với các đường thẳng ${d_1}$, ${d_2}$, ${d_3}$.
	 \loigiai{
		\begin{center}
			\begin{tikzpicture}[scale=1,line cap=round,line join=round,font=\footnotesize,>=stealth]
					\def\dai{4}  \def\rong{2} \def\goc{50}
					\path (0,0) coordinate (O)++(0:\dai) coordinate  (O1)
						++(\goc:\rong) coordinate (O2)++(180:\dai)  coordinate (O3)
						($(O)+(40:1.2)$) coordinate (P)++(0:0.5) coordinate (A)++(0:0.8) coordinate (A1)++(0:1.5)  coordinate (B)++(0:0.5) coordinate (Q)
						($(P)+(90:1.5)$) coordinate (P')++(0:1) coordinate (C)++(0:3.2) coordinate (Q'); 
					\coordinate (C1) at ($(C)!-0.3!(A1)$); 
					 \coordinate (C2) at ($(C)!2!(A1)$);
					\coordinate (C3)  at (intersection of O--O1 and C--C2);
					\draw (O)--(O1)--(O2)--(O3)--cycle (P)--node[pos=0.7,below]{$d_1$} 
						(Q) 
						 (P')--node[pos=0.7,above]{$d_3$}
						(Q') 
						(C1)--(A1) 
						(C2)--node[pos=0.4,left]{$d_2$} 
						 (C3);
					 \draw [dashed] (A1)--(C3);
					\draw  pic[draw,angle radius=7mm]{angle=O1--O--O3};
					\draw (O)  node[shift={(25:0.5)}]{$P$};
					\foreach  \t/\g in {A/90,B/90,C/50}
					 \draw[fill] (\t) circle (1pt)
						node[shift={(\g:7pt)}]{$\t$}; 
				\end{tikzpicture}
		\end{center}
		Đường thẳng ${d_1}$ chứa  hai điểm $A$, $B$ thuộc $( P)$, vậy ${d_1} \subset ( P)$.\\
		Đường thẳng ${d_2}$ không nằm trong $( P)$, vì có chứa điểm $C$ không thuộc $( P)$. Mặt khác, ${d_2}$ lại có điểm $A$ chung với $( P)$, suy ra ${d_2}$ cắt $( P)$  tại $A$.\\
		Đường thẳng ${d_3}$ không nằm trong $( P)$ và song song với đường thẳng ${d_1}$ nằm trong  $( P)$, suy ra ${d_3}\parallel ( P)$.
	} 
\end{ex}

% nguồn: Ok_TuanLuan_Toan11_Lan2
\begin{ex}%[1H4H3-1]%[ma:7e9ecc37]
	Cho hai hình bình hành $ABCD$ và $ABMN$ không đồng phẳng, xác định vị trí  tương đối của mặt phẳng $( ABMN)$ với các đường thẳng $CD$, $BD$ và $BN$.
	 \loigiai{
		\begin{center}
			\begin{tikzpicture}[scale=1,line cap=round,line join=round,font=\footnotesize,>=stealth]
					\def\dai{4} \def\rong{1.6}  \def\goc{-110}
					\path (0,0) coordinate  (A)++(0:\dai) coordinate (B)
						++(\goc:\rong)  coordinate (C)++(-180:\dai) coordinate (D)
						($(A)+(110:\rong)$) coordinate (N)++(0:\dai)  coordinate (M);
					\draw (A) --(B) --(C) --(D) --cycle 
						 (A)--(N) --(M) --(B)
						 (N)--(B)--(D);
					\foreach \t/\g  in {A/180,B/0,C/0,D/180,M/0,N/180}
					 \draw[fill] (\t) circle (1.0pt)
						 node[shift={(\g:9pt)}]{$\t$};
				\end{tikzpicture}
		\end{center}
		Nếu $CD$ có điểm chung $O$ với mặt phẳng $( ABMN)$ thì $O$ thuộc giao tuyến $AB$ của hai mặt phẳng $( ABCD)$ và $( ABMN)$, suy ra $AB$ cắt $CD$ (Mâu thuẫn với giả thiết $ABCD$ là hình bình  hành). \\
		Vậy $CD\parallel ( ABMN)$.\\ 
		Đường thẳng $BD$ có  một điểm chung duy nhất $B$ với $( ABMN) \Rightarrow BD$ cắt $( ABMN)$ tại $B$.\\
		Đường thẳng $BN$ có hai điểm chung với $( ABMN)$  nên $BN \subset ( ABMN)$.
	 }
\end{ex}

% nguồn: Ok_TuanLuan_Toan11_Lan2
\begin{ex}%[1H4H3-1]%[ma:ce626905]
	Cho hai hình bình hành $ABCD$ và  $ABMN$ không đồng phẳng. Xác định vị trí tương đối của mặt phẳng $( ABCD)$ lần lượt với các đường thẳng $CD$, $BD$ và $BN$.
	 \loigiai{
		\begin{center}
			\begin{tikzpicture}[scale=1,line cap=round,line join=round,font=\footnotesize,>=stealth]
					\def\dai{4} \def\rong{1.6} \def\goc{-110} 
					\path (0,0)  coordinate (A)++(0:\dai) coordinate (B)
						++(\goc:\rong) coordinate (C)++(-180:\dai) coordinate  (D)
						($(A)+(110:\rong)$) coordinate  (N)++(0:\dai) coordinate (M);
					\draw (A) node[left]{$A$}--(B) node[right]{$B$}--(C) node[right]{$C$}--(D)  node[left]{$D$}--cycle
						(A)--(N) node[left]{$N$}--(M)  node[right]{$M$}--(B)
						(N)--(B)--(D); 
				\end{tikzpicture}
		\end{center}
		Nếu $CD$ có điểm chung $O$ với $( ABMN)$ thì $O$ thuộc giao tuyến  $AB$ của hai mp $( ABCD)$ và $( ABMN)$ suy ra $CD$ cắt $AB$ (mâu thuẫn với giả thiết $ABCD$ là hình bình hành).\\
		$BD$ có một điểm chung duy nhất $B$ với $( ABMN)$ suy ra  $BD$ cắt $( ABMN)$ tại $B$.\\
		$BN$ có hai điểm chung $B$ và $N$ với $( ABMN)$ suy ra  $BN \subset ( ABMN)$.
	} 
\end{ex}
